\documentclass[12pt]{article}
\usepackage[utf8]{inputenc}
\usepackage[T1]{fontenc}
\usepackage[margin=1in]{geometry}
\usepackage{setspace}
\usepackage{xcolor}
\usepackage[colorlinks=true, urlcolor=blue!60!black, linkcolor=blue!60!black]{hyperref}
\setlength{\parindent}{2em}
\setlength{\parskip}{1.25em}
\renewcommand{\baselinestretch}{1.0}
\begin{document}
\begin{center}
{\LARGE \textit{Nevada Department of Human Resources v. Hibbs}, 538 U.S. 721 (2003)\par}
\vspace{0.8em}
{\large Case Brief\par}
\end{center}
\vspace{1em}

\section{Memory Jogger}
In a 6--3 decision authored by Chief Justice Rehnquist, the Supreme Court held that Congress validly abrogated state Eleventh Amendment sovereign immunity under Section 5 of the Fourteenth Amendment when it enacted the Family and Medical Leave Act's family-care leave provision because gender classifications receive heightened scrutiny and the gender-neutral 12-week leave floor is a congruent and proportional prophylactic remedy against pervasive state sex-role stereotyping.

\section{Detailed Case Facts}
Enacted in 1993, the Family and Medical Leave Act (FMLA), 29 U.S.C. \S\S~2601--2654, entitles eligible employees to take up to 12 workweeks of unpaid, job-protected leave annually for specified family and medical reasons, including the onset of a ``serious health condition'' in an employee's spouse, child, or parent under the family-care provision, 29 U.S.C. \S~2612(a)(1)(C) (\href{https://supreme.justia.com/cases/federal/us/538/721/}{538 U.S. at 724}). The Act prohibits employers from interfering with, restraining, or denying the exercise of FMLA rights (\S~2615(a)(1)) and creates a private right of action for equitable relief and money damages ``against any employer (including a public agency) in any Federal or State court of competent jurisdiction'' (\S~2617(a)(2)), defining ``public agency'' to include state governments and state agencies (\S\S~203(x), 2611(4)(A)(iii)) (\href{https://supreme.justia.com/cases/federal/us/538/721/}{538 U.S. at 724--726}).

Respondent William Hibbs was employed in the Welfare Division of the Nevada Department of Human Resources (\href{https://supreme.justia.com/cases/federal/us/538/721/}{538 U.S. at 725}). In April and May 1997, Hibbs requested leave under the FMLA to care for his wife, who was recovering from a serious automobile accident and neck surgery (\href{https://supreme.justia.com/cases/federal/us/538/721/}{538 U.S. at 725}; \href{https://www.oyez.org/cases/2002/01-1368}{Oyez Case Summary}). The Department approved his request for the full 12 weeks of FMLA leave and authorized him to take the leave intermittently as needed between May and December 1997. Hibbs used intermittent leave until August 5, 1997, after which he did not return to work. In October 1997, the Department notified Hibbs that he had exhausted his 12 weeks of FMLA leave, that no additional leave would be granted, and that he had to report back to work by November 12, 1997. When Hibbs failed to report on that date, the Department terminated his employment (\href{https://supreme.justia.com/cases/federal/us/538/721/}{538 U.S. at 725}).

\section{Procedural History}
Hibbs sued the Nevada Department of Human Resources and two of its officers in the U.S. District Court for the District of Nevada, seeking money damages as well as injunctive and declaratory relief for, among other claims, violations of 29 U.S.C. \S~2612(a)(1)(C) (\href{https://supreme.justia.com/cases/federal/us/538/721/}{538 U.S. at 725}). The District Court granted summary judgment to the state defendants, holding that the Eleventh Amendment barred Hibbs's FMLA damages claim against the State and that his Fourteenth Amendment rights had not been violated.

Hibbs appealed to the U.S. Court of Appeals for the Ninth Circuit, and the United States intervened pursuant to 28 U.S.C. \S~2403 to defend the constitutionality of the FMLA's abrogation of state sovereign immunity (\href{https://supreme.justia.com/cases/federal/us/538/721/}{538 U.S. at 725}). The Ninth Circuit reversed, 273 F.3d 844 (9th Cir. 2001), concluding that Congress validly abrogated state sovereign immunity under Section 5 of the Fourteenth Amendment. Because the Ninth Circuit's decision created a circuit split with the Fifth Circuit's contrary ruling in \textit{Kazmier v. Widmann}, 225 F.3d 519 (5th Cir. 2000), and other circuits, the Supreme Court granted certiorari, 536 U.S. 938 (2002) (\href{https://supreme.justia.com/cases/federal/us/538/721/}{538 U.S. at 725--726}). The case (No. 01-1368) was argued on January 15, 2003, and decided on May 27, 2003 (\href{https://www.courtlistener.com/opinion/127926/nevada-department-of-human-resources-v-hibbs/}{CourtListener Opinion}; \href{https://www.oyez.org/cases/2002/01-1368}{Oyez Case Summary}).

\section{Judicial Votes}
\begin{itemize}
\item \textbf{Vote:} 6--3, affirming the judgment of the Ninth Circuit (\href{https://supreme.justia.com/cases/federal/us/538/721/}{538 U.S. at 723}).
\item \textbf{Majority (5):} Chief Justice Rehnquist delivered the opinion of the Court, joined by Justices O'Connor, Souter, Ginsburg, and Breyer.
\item \textbf{Concurrence (3):} Justice Souter filed a concurring opinion, joined by Justices Ginsburg and Breyer.
\item \textbf{Concurrence in the Judgment (1):} Justice Stevens filed an opinion concurring in the judgment only.
\item \textbf{Dissents (3):} Justice Scalia filed a dissenting opinion; Justice Kennedy filed a dissenting opinion, joined by Justices Scalia and Thomas.
\end{itemize}

\section{Holding}
\textbf{Judgment:} Affirmed. State employees may recover money damages in federal court against a nonconsenting State for the State's failure to comply with the family-care leave provision of the FMLA, 29 U.S.C. \S~2612(a)(1)(C) (\href{https://supreme.justia.com/cases/federal/us/538/721/}{538 U.S. at 725, 740}).

\textbf{Rule Announced:} Congress may abrogate States' Eleventh Amendment sovereign immunity in federal court when it (1) makes its intent to abrogate unmistakably clear in the statutory text, and (2) acts pursuant to a valid exercise of its enforcement power under Section 5 of the Fourteenth Amendment (\href{https://supreme.justia.com/cases/federal/us/538/721/}{538 U.S. at 726}). Under \textit{City of Boerne v. Flores}, 521 U.S. 507 (1997), prophylactic Section 5 legislation must exhibit ``congruence and proportionality between the injury to be prevented or remedied and the means adopted to that end'' (\href{https://supreme.justia.com/cases/federal/us/538/721/}{538 U.S. at 728}). Crucially, when Congress targets state discrimination involving a classification subject to heightened constitutional scrutiny---such as sex discrimination, which triggers intermediate scrutiny under \textit{Craig v. Boren} and \textit{United States v. Virginia}---it is significantly easier for Congress to demonstrate a pattern of unconstitutional state conduct justifying prophylactic Section 5 remedies than when Congress targets nonsuspect classifications such as age or disability (\href{https://supreme.justia.com/cases/federal/us/538/721/}{538 U.S. at 735--736}).

\textbf{Holding vs.\ Dicta:} The Court's binding holding applies specifically to the FMLA's \textit{family-care} leave provision, 29 U.S.C. \S~2612(a)(1)(C), holding that Congress satisfied the clear-statement rule and validly exercised its Section 5 power to abrogate state sovereign immunity for that provision (\href{https://supreme.justia.com/cases/federal/us/538/721/}{538 U.S. at 725, 740}). The opinion did \textit{not} hold that the FMLA's \textit{self-care} provision (\S~2612(a)(1)(D), covering an employee's own serious health condition) validly abrogates state immunity---a distinction the Supreme Court later made dispositive in \textit{Coleman v. Court of Appeals of Maryland}, 566 U.S. 30 (2012), which held that suits against States under the self-care provision are barred by sovereign immunity (\href{https://supreme.justia.com/cases/federal/us/566/30/}{\textit{Coleman}, 566 U.S. 30 (2012)}). Likewise, the majority's footnote 12 discussion of how existing Title VII and Equal Protection rules would independently bar a State from offering 24 weeks of leave to women and only 12 weeks to men (\href{https://supreme.justia.com/cases/federal/us/538/721/}{538 U.S. at 739 n. 12}) is explanatory commentary rather than a separate holding.

\section{Analysis of Opinions}

\subsection{Majority Opinion (Chief Justice Rehnquist)}
Chief Justice Rehnquist---surprising many observers given his leadership of the Court's earlier federalism revival---wrote for a five-Justice majority upholding the FMLA's abrogation of state sovereign immunity (\href{https://supreme.justia.com/cases/federal/us/538/721/}{538 U.S. at 724--740}; \href{https://www.courtlistener.com/opinion/127926/nevada-department-of-human-resources-v-hibbs/}{CourtListener Opinion}). First, he easily found that the FMLA satisfied the clear-statement rule of \textit{Dellmuth v. Muth} because its definition of ``public agency'' mirrored the language upheld as clear in \textit{Kimel v. Florida Board of Regents}, 528 U.S. 62 (2000). Because \textit{Seminole Tribe of Florida v. Florida}, 517 U.S. 44 (1996), forecloses abrogation under Article I's Commerce Clause, the case turned on Section 5 of the Fourteenth Amendment under \textit{Fitzpatrick v. Bitzer}, 427 U.S. 445 (1976), and \textit{City of Boerne v. Flores}, 521 U.S. 507 (1997) (\href{https://supreme.justia.com/cases/federal/us/538/721/}{538 U.S. at 726--728}).

Second, Chief Justice Rehnquist surveyed the long history of state-sanctioned sex discrimination in employment---from \textit{Bradwell v. Illinois} (1873) and \textit{Muller v. Oregon} (1908) to \textit{Goesaert v. Cleary} (1948) and \textit{Hoyt v. Florida} (1961)---and examined the congressional record behind the FMLA (\href{https://supreme.justia.com/cases/federal/us/538/721/}{538 U.S. at 729--734}). That record included 1989 and 1990 Bureau of Labor Statistics surveys showing a widening gap between maternity and paternity leave (37\% versus 18\% in 1990), a 50-state study showing public-sector leave policies mirrored private-sector disparities, testimony that fathers faced ``notoriously discriminatory treatment'' when requesting parental leave, data showing 15 States offered women up to a year of extended maternity leave (far beyond the 4-to-8-week medical recovery period) while only four States offered men parallel leave, and evidence that discretionary leave policies invited unequal application (\href{https://supreme.justia.com/cases/federal/us/538/721/}{538 U.S. at 730--732}).

Third, Chief Justice Rehnquist distinguished \textit{Kimel} (ADEA) and \textit{Board of Trustees of the University of Alabama v. Garrett}, 531 U.S. 356 (2001) (ADA Title I). Age and disability classifications receive only rational-basis review, so Congress had to document a widespread pattern of \textit{irrational} state discrimination. Sex classifications, by contrast, trigger heightened (intermediate) scrutiny under \textit{Craig v. Boren} and \textit{United States v. Virginia}, making it ``easier for Congress to show a pattern of state constitutional violations'' (\href{https://supreme.justia.com/cases/federal/us/538/721/}{538 U.S. at 735--736}; \href{https://www.law.cornell.edu/constitution-conan/amendment-14/section-5/modern-doctrine}{LII Constitution Annotated}). Mutually reinforcing stereotypes---that women are primary caregivers and men lack domestic responsibilities---created a self-fulfilling cycle of subtle discrimination that Title VII and the Pregnancy Discrimination Act had failed to eradicate.

Finally, the majority held that the FMLA's 12-week gender-neutral leave floor was congruent and proportional. A statute merely prohibiting sex discrimination in leave policies would allow States to offer no family leave at all, disproportionately excluding working women (\href{https://supreme.justia.com/cases/federal/us/538/721/}{538 U.S. at 738}). Moreover, unlike the sweeping statutes invalidated in \textit{Boerne}, \textit{Kimel}, and \textit{Garrett}, the FMLA was ``narrowly targeted at the faultline between work and family'' and contained substantial built-in limits: leave is unpaid, requires one year and 1,250 hours of prior service, excludes elected officials and high-ranking policymakers, requires advance notice and medical certification, and caps damages at actual monetary losses with a two-year (three years if willful) limitations period (\href{https://supreme.justia.com/cases/federal/us/538/721/}{538 U.S. at 738--740}).

\subsection{Concurring Opinion (Justice Souter)}
Justice Souter, joined by Justices Ginsburg and Breyer, wrote a one-paragraph concurrence joining the Chief Justice's opinion while preserving their broader views of federal power (\href{https://supreme.justia.com/cases/federal/us/538/721/}{538 U.S. at 740}). He noted that even under the Court's restrictive Section 5 framework in \textit{Garrett}, \textit{Kimel}, and \textit{Florida Prepaid}, the FMLA was ``undoubtedly valid legislation,'' and that the same conclusion followed \textit{a fortiori} from the dissenters' more expansive understanding of Section 5 and their rejection of \textit{Seminole Tribe}'s state sovereign immunity doctrine.

\subsection{Opinion Concurring in the Judgment (Justice Stevens)}
Justice Stevens concurred in the judgment only, refusing to join the Section 5 analysis (\href{https://supreme.justia.com/cases/federal/us/538/721/}{538 U.S. at 740--741}). Expressing doubt that the FMLA was truly ``needed to secure the guarantees of the Fourteenth Amendment,'' he grounded his vote instead in his longstanding critique of \textit{Hans v. Louisiana} and \textit{Seminole Tribe}. Because Hibbs was a citizen of Nevada, the plain text of the Eleventh Amendment did not bar his suit; Nevada's immunity defense rested solely on judge-made common law, which Congress could freely abrogate under its Article I Commerce Clause power (\href{https://supreme.justia.com/cases/federal/us/538/721/}{538 U.S. at 741}).

\subsection{Dissenting Opinion (Justice Scalia)}
Justice Scalia joined Justice Kennedy's dissent and added a sharp federalism objection (\href{https://supreme.justia.com/cases/federal/us/538/721/}{538 U.S. at 741--743}). Relying on the \textit{Civil Rights Cases}, 109 U.S. 3 (1883), he argued that ``[t]here is no guilt by association'' under Section 5: Congress cannot abrogate one State's sovereign immunity because of constitutional violations committed by other States. While prophylaxis may extend a remedy beyond the exact constitutional violation, it may not extend ``beyond the violator'' (\href{https://supreme.justia.com/cases/federal/us/538/721/}{538 U.S. at 742}). Even if \S~2612(a)(1)(C) survived a facial challenge, Justice Scalia insisted that Nevada remained entitled on an as-applied challenge to demand proof that Nevada itself had violated the Fourteenth Amendment (\href{https://supreme.justia.com/cases/federal/us/538/721/}{538 U.S. at 743}).

\subsection{Dissenting Opinion (Justice Kennedy)}
Justice Kennedy, joined by Justices Scalia and Thomas, argued that the majority gutted \textit{Boerne}'s congruence-and-proportionality test and permitted Congress to impose a substantive federal entitlement program on state treasuries (\href{https://supreme.justia.com/cases/federal/us/538/721/}{538 U.S. at 744--759}). Scrutinizing the legislative record, Justice Kennedy contended that the BLS statistics surveyed only private employers; the 1986 Yale Bush Center and Washington Council of Lawyers testimony concerned \textit{parenting} leave rather than family-care leave under \S~2612(a)(1)(C); the discretionary-leave report examined federal agencies, not States; and 30 States, the District of Columbia, and Puerto Rico had already adopted family-care leave laws before 1993---including Nevada, which allowed up to a year of gender-neutral unpaid leave, 10 days of paid family sick leave, and a catastrophic-leave bank (\href{https://supreme.justia.com/cases/federal/us/538/721/}{538 U.S. at 746--755}). Stripped of pregnancy-disability schemes valid under \textit{Geduldig v. Aiello}, 417 U.S. 484 (1974), the record showed at most three States (Massachusetts, Kansas, and Tennessee) with female-only parenting leave. Heightened scrutiny, Justice Kennedy argued, did not relieve Congress of its burden to document a pattern of unconstitutional state conduct. Finally, because the FMLA mandated a 12-week floor yet did not even bar a State from giving women 24 weeks and men 12 weeks, it operated as a substantive benefit program rather than a remedial statute; plaintiffs still had adequate recourse through Commerce Clause enforcement by the United States and \textit{Ex parte Young} injunctive suits (\href{https://supreme.justia.com/cases/federal/us/538/721/}{538 U.S. at 755--759}).

\section{Examples: Future Applications}
\begin{enumerate}
\item \textbf{Same-Side (Valid Abrogation): Federal Protections Against Race-Based Disparate Impact in State Employment.} Suppose Congress re-enacts or strengthens Title VII's disparate-impact provisions authorizing state employees to sue States for money damages over facially neutral state hiring tests that disproportionately exclude racial minorities without business necessity. Under \textit{Hibbs}, this falls on the same side and is a valid exercise of Section 5 power. Because racial classifications trigger strict scrutiny---even higher than the intermediate scrutiny for sex in \textit{Hibbs}---Congress faces a lower evidentiary hurdle to document a pattern of unconstitutional state race discrimination, and targeted prophylactic remedies aimed at subtle workplace bias are congruent and proportional (\href{https://openscholarship.wustl.edu/cgi/viewcontent.cgi?referer=&httpsredir=1&article=1395&context=law_journal_law_policy}{Chemerinsky Analysis}).

\item \textbf{Same-Side (Valid Abrogation): ADA Title II Claims Implicating Fundamental Constitutional Rights.} Suppose a paraplegic litigant sues a State for money damages under Title II of the Americans with Disabilities Act because a county courthouse lacks an elevator, forcing him to crawl up stairs to attend court proceedings. Under \textit{Hibbs}'s heightened-scrutiny framework---as directly applied by the Supreme Court in \textit{Tennessee v. Lane}, 541 U.S. 509 (2004)---this falls on the same side. Even though disability classifications receive rational-basis review, denying physical access to the courts implicates a fundamental due-process right subject to heightened scrutiny, making Title II's reasonable-modification requirement a congruent and proportional Section 5 remedy (\href{https://supreme.justia.com/cases/federal/us/541/509/}{\textit{Tennessee v. Lane}, 541 U.S. 509 (2004)}).

\item \textbf{Opposite-Side (Invalid Abrogation): FMLA Self-Care Leave Suits Against States.} Suppose a male state court employee is fired after requesting leave for his own serious medical condition under the FMLA's self-care provision, 29 U.S.C. \S~2612(a)(1)(D), and sues the State for money damages. This falls on the opposite side and is barred by state sovereign immunity, as the Supreme Court held in \textit{Coleman v. Court of Appeals of Maryland}, 566 U.S. 30 (2012) (\href{https://supreme.justia.com/cases/federal/us/566/30/}{\textit{Coleman}, 566 U.S. 30 (2012)}). Unlike the family-care provision in \textit{Hibbs}, Congress compiled no legislative record showing that States administered personal sick-leave policies in a sex-discriminatory manner.

\item \textbf{Opposite-Side (Invalid Abrogation): Federal Mandate Requiring Paid Vacation Leave to Combat Socioeconomic Inequality.} Suppose Congress enacts a statute requiring all public and private employers to provide four weeks of paid annual vacation and authorizes state employees to sue States for money damages, citing general wealth disparities. This falls squarely on the opposite side. Wealth is not a suspect classification under the Equal Protection Clause, so rational-basis review applies (as in \textit{Kimel} and \textit{Garrett}), and imposing an across-the-board paid benefit with no documented pattern of unconstitutional state conduct is a substantive entitlement program exceeding Section 5.

\item \textbf{Fence-Sitter (Close Case): Federal Reasonable-Accommodation Mandate for Pregnant State Workers (e.g., the Pregnant Workers Fairness Act).} Suppose a pregnant state employee sues a State for money damages under a federal statute requiring employers to provide reasonable accommodations for known limitations related to pregnancy, childbirth, or related medical conditions. On one hand, under \textit{Hibbs}, pregnancy and workplace caregiving lie directly along ``the faultline between work and family'' where sex-based stereotyping is strongest, and the statute contains targeted limits (such as an ``undue hardship'' defense) analogous to the FMLA's limits. On the other hand, under \textit{Geduldig v. Aiello} (cited by Justice Kennedy's dissent in \textit{Hibbs} and echoed in \textit{Coleman}), pregnancy-based distinctions are not automatically treated as facial sex classifications under the Equal Protection Clause, giving States a strong argument that Congress must compile a rigorous state-specific record of intentional sex discrimination to abrogate sovereign immunity.
\end{enumerate}

\section{Critique}
\textit{Hibbs} sits at the intersection of federalism, equal protection, and separation of powers, drawing sharp commentary from both conservative and progressive constitutional scholars. From an \textbf{originalist and federalism-centered perspective}, critics argued that \textit{Hibbs} inconsistently relaxed \textit{City of Boerne}'s discipline. Echoing Justices Kennedy and Scalia, conservative commentators contended that Chief Justice Rehnquist relied on thin, private-sector, and parenting-leave evidence to validate a nationwide substantive labor mandate against all 50 States---including States like Nevada that had progressive, gender-neutral leave policies (\href{https://supreme.justia.com/cases/federal/us/538/721/}{538 U.S. at 741--755}; \href{https://www.cato.org/supreme-court-review/2003-2004/tennessee-v-lane-how-illegitimate-power-negated-non-existent-immunity}{Cato Supreme Court Review}). Some libertarian and textualist scholars also observed a deeper irony: the entire Section 5 gymnastics in \textit{Hibbs} was only necessary because the Court's nontextual sovereign-immunity doctrine in \textit{Hans} and \textit{Seminole Tribe} barred Congress from using its plain Article I Commerce Clause power to authorize damages suits against States (\href{https://www.cato.org/supreme-court-review/2003-2004/tennessee-v-lane-how-illegitimate-power-negated-non-existent-immunity}{Cato Supreme Court Review}).

From a \textbf{progressive and doctrinal perspective}, scholars celebrated \textit{Hibbs} as a landmark victory for sex equality and civil rights while critiquing the underlying \textit{Boerne} framework. Robert Post and Reva Siegel argued that \textit{Hibbs} implicitly vindicated a ``policentric'' model of constitutional interpretation in which Congress, responsive to social movements and democratic mobilization, can enforce evolving understandings of equal protection at the work-family boundary rather than merely mimicking judicial doctrine (\href{https://law.yale.edu/sites/default/files/documents/pdf/Faculty/Siegel_LegislativeConstitutionalismAndSectionFivePower.pdf}{Post \& Siegel, Yale Law}). Samuel Bagenstos highlighted how \textit{Hibbs} endorsed ``feminist universalism''---recognizing that universal, gender-neutral family-leave floors dismantle sex-role stereotypes more effectively than formal antidiscrimination bans alone (\href{https://repository.law.umich.edu/cgi/viewcontent.cgi?article=1182&context=book_chapters}{Bagenstos, Univ.\ of Michigan Law}). At the same time, Erwin Chemerinsky and other progressive scholars noted that tying Section 5 power to judicial tiers of scrutiny creates an arbitrary constitutional hierarchy that shields sex- and race-based legislation while leaving older and disabled state workers unprotected against damages suits (\href{https://openscholarship.wustl.edu/cgi/viewcontent.cgi?referer=&httpsredir=1&article=1395&context=law_journal_law_policy}{Chemerinsky, WashU Law}).

Underlying these debates is a fundamental clash of \textbf{competing constitutional values}: (1) \textit{state sovereign dignity, fiscal autonomy, and judicial supremacy} in defining constitutional rights under \textit{Boerne} and \textit{Seminole Tribe} versus (2) \textit{congressional enforcement authority and substantive gender equality} under the Reconstruction Amendments' mandate to eradicate entrenched state-fostered discrimination.

\section{Key Quotations}

\begin{quote}
``Stereotypes about women's domestic roles are reinforced by parallel stereotypes presuming a lack of domestic responsibilities for men. Because employers continued to regard the family as the woman's domain, they often denied men similar accommodations or discouraged them from taking leave. These mutually reinforcing stereotypes created a self-fulfilling cycle of discrimination that forced women to continue to assume the role of primary family caregiver, and fostered employers' stereotypical views about women's commitment to work and their value as employees.''
\end{quote}
\noindent --- Chief Justice Rehnquist, Opinion of the Court, \href{https://supreme.justia.com/cases/federal/us/538/721/}{538 U.S. 721, 736 (2003)} (also on \href{https://www.courtlistener.com/opinion/127926/nevada-department-of-human-resources-v-hibbs/}{CourtListener}).

\begin{quote}
``By setting a minimum standard of family leave for all eligible employees, irrespective of gender, the FMLA attacks the formerly state-sanctioned stereotype that only women are responsible for family care giving, thereby reducing employers' incentives to engage in discrimination by basing hiring and promotion decisions on stereotypes.''
\end{quote}
\noindent --- Chief Justice Rehnquist, Opinion of the Court, \href{https://supreme.justia.com/cases/federal/us/538/721/}{538 U.S. 721, 737 (2003)}.

\begin{quote}
``Unlike the statutes at issue in \textit{City of Boerne}, \textit{Kimel}, and \textit{Garrett}, which applied broadly to every aspect of state employers' operations, the FMLA is narrowly targeted at the faultline between work and family---precisely where sex-based overgeneralization has been and remains strongest---and affects only one aspect of the employment relationship.''
\end{quote}
\noindent --- Chief Justice Rehnquist, Opinion of the Court, \href{https://supreme.justia.com/cases/federal/us/538/721/}{538 U.S. 721, 738 (2003)}.

\begin{quote}
``The constitutional violation that is a prerequisite to `prophylactic' congressional action to `enforce' the Fourteenth Amendment is a violation by the State against which the enforcement action is taken. There is no guilt by association, enabling the sovereignty of one State to be abridged under \S~5 of the Fourteenth Amendment because of violations by another State, or by most other States, or even by 49 other States.''
\end{quote}
\noindent --- Justice Scalia, Dissenting Opinion, \href{https://supreme.justia.com/cases/federal/us/538/721/}{538 U.S. 721, 741--742 (2003)}.

\begin{quote}
``The Court is unable to show that States have engaged in a pattern of unlawful conduct which warrants the remedy of opening state treasuries to private suits. The inability to adduce evidence of alleged discrimination, coupled with the inescapable fact that the federal scheme is not a remedy but a benefit program, demonstrates the lack of the requisite link between any problem Congress has identified and the program it mandated.''
\end{quote}
\noindent --- Justice Kennedy, Dissenting Opinion, \href{https://supreme.justia.com/cases/federal/us/538/721/}{538 U.S. 721, 745 (2003)}.

\end{document}
